\documentclass{article}

\PassOptionsToPackage{numbers}{natbib}
\usepackage[final]{formatting/neurips_2024}
\usepackage{enumitem}
\usepackage{tcolorbox}
\tcbuselibrary{skins,breakable}
\usepackage[utf8]{inputenc}
\usepackage[T1]{fontenc}
\usepackage{hyperref}
\usepackage{url}
\usepackage{booktabs}
\usepackage{amsfonts}
\usepackage{nicefrac}
\usepackage{microtype}
\usepackage{xcolor}
\usepackage{amsmath}
\usepackage{amsthm}
\usepackage{algorithm}
\usepackage{algpseudocode}
\usepackage{titlesec}
\usepackage{graphicx}
\usepackage{mathtools} % provides \coloneqq
\usepackage{multirow}

\setlist{nosep,leftmargin=*}
\setlist[enumerate]{label=(\arabic*)}
\setlength{\abovedisplayskip}{4pt}
\setlength{\belowdisplayskip}{4pt}
\setlength{\abovedisplayshortskip}{2pt}
\setlength{\belowdisplayshortskip}{2pt}
\setlength{\textfloatsep}{8pt}
\setlength{\floatsep}{6pt}
\setlength{\intextsep}{6pt}
\titlespacing*{\section}{0pt}{1.2ex}{0.6ex}
\titlespacing*{\subsection}{0pt}{1ex}{0.4ex}

\newtheorem{lemma}{Lemma}
\newtheorem{remark}{Remark}
\newtheorem{proposition}{Proposition}

\usepackage{color-edits}
\addauthor[Ellen]{ev}{red}
\addauthor[SI]{si}{teal}

\title{Caching}

\begin{document}
\maketitle

\section{Experiments}

\subsection{Setup}

Dataset: 3000 synthetic traces (1000 per family LRU/LFU/ARC), each of length
$T{=}16{,}000$ over a universe of $U{=}512$ items, cache capacity $k{=}32$.
Each family is filtered so its target algorithm beats the other two by at
least $0.04$ absolute hit-rate gap. Split is per-trace-row 80/10/10
(train/val/test), seed=0, preventing leakage across timesteps of the same
trace.

Model: causal transformer with a dual-attention block (one head keyed on the
full cache, one head keyed on a sliding window of the recent request
sequence), $d_{\text{model}}{=}128$, $d_{\text{ff}}{=}256$, 2 layers, 2 heads,
context window $w{=}1024$, $465{,}665$ parameters. Trained for 6 epochs under
teacher forcing: at every Belady full-cache-miss event the model sees
Belady's own cache state and predicts the slot to evict, supervised by
Belady's choice with cross-entropy loss. Optimisation: AdamW, lr $10^{-3}$,
effective batch 1024 (4$\times$ A6000 GPUs via DDP, 256 per rank).

\subsection{Training curves}

Training and validation loss/accuracy per epoch are plotted in
\texttt{learned\_eviction/runs/default/training\_curves.png}. Train loss
descends rapidly from chance ($\log 32 \approx 3.47$) around step 800, plateaus
near 1.8 through most of epoch 0, breaks below 1.5 mid-epoch 2, and then
bounces in the $1.5\text{--}1.8$ band through epochs 3--5. Per-epoch
validation loss decreases monotonically (epochs 0--4) from $1.7825$ to
$1.7447$ with an uptick at epoch 5 ($1.7497$). Validation accuracy inches up
every epoch (51.9\% $\to$ 52.5\%). Best checkpoint (lowest val loss) is
epoch 4.

\subsection{Teacher-forced test evaluation}

The model is evaluated on the held-out test set (1.28M eviction events
across 300 traces, 100 per family) with Belady's cache state as input
(identical to training conditions). Item-level agreement with LRU/LFU/ARC is
measured at \emph{shared} eviction timesteps: for each Belady eviction event
at time $t$, if algorithm $X$ also evicts at $t$ (on its own running cache),
compare the item that the transformer evicts (from Belady's cache) to the
item $X$ evicts. Coverage reports the fraction of Belady events where $X$
also performs an eviction; agreement is conditional on coverage.

\paragraph{Overall.} Test loss $1.7374$, slot-accuracy (agreement with
Belady) $52.46\%$. This matches epoch-5 validation (52.48\%), indicating no
train/test generalisation gap.

\paragraph{Per-dominating-family accuracy.} Grouping test examples by the
dominant algorithm of their source trace (determined from the recorded
hit-rate scores in \texttt{*\_scores.json}, which by construction matches
the source family) reveals a stark asymmetry:

\begin{center}
\begin{tabular}{lrr}
\toprule
Dominant family & \#examples & Slot-accuracy \\
\midrule
LRU-dominant & 306,417 & 13.1\% \\
LFU-dominant & 528,821 & 61.4\% \\
ARC-dominant & 445,781 & 68.9\% \\
\bottomrule
\end{tabular}
\end{center}

\paragraph{Per-dominating-family agreement with each classical algorithm.}

\begin{center}
\begin{tabular}{llrr}
\toprule
Group & Algo & Coverage & Agreement \\
\midrule
\multirow{3}{*}{All} & LRU & $100.0\%$ & $\phantom{0}0.3\%$ \\
                    & LFU & $\phantom{0}92.8\%$  & $46.0\%$ \\
                    & ARC & $\phantom{0}99.7\%$  & $\phantom{0}9.8\%$ \\
\midrule
\multirow{3}{*}{LRU-dom} & LRU & $100.0\%$ & $\phantom{0}1.2\%$ \\
                         & LFU & $\phantom{0}94.1\%$  & $\phantom{0}7.3\%$ \\
                         & ARC & $\phantom{0}99.9\%$  & $\phantom{0}1.7\%$ \\
\midrule
\multirow{3}{*}{LFU-dom} & LRU & $100.0\%$ & $\phantom{0}0.0\%$ \\
                         & LFU & $\phantom{0}90.3\%$  & $64.9\%$ \\
                         & ARC & $\phantom{0}99.4\%$  & $\phantom{0}5.5\%$ \\
\midrule
\multirow{3}{*}{ARC-dom} & LRU & $100.0\%$ & $\phantom{0}0.0\%$ \\
                         & LFU & $\phantom{0}94.8\%$  & $51.0\%$ \\
                         & ARC & $\phantom{0}99.8\%$  & $20.5\%$ \\
\bottomrule
\end{tabular}
\end{center}

\subsection{Interpretation}

\paragraph{The model has learned an LFU-flavoured policy.} LFU agreement
dominates every row: $46\%$ overall, $65\%$ on LFU-dominant traces, and
still $51\%$ even on ARC-dominant traces. Agreement with LRU is essentially
zero ($0.3\%$ overall), and agreement with ARC is modest ($\sim\!10\%$).
Whatever heuristic the transformer has internalised, it looks much more like
``evict the least-frequently-accessed item'' than like recency tracking or
ARC's T1/T2 logic.

\paragraph{The headline 52\% accuracy hides a sharp per-family split.} The
model is near-optimal on ARC traces ($69\%$ agreement with Belady),
competent on LFU traces ($61\%$), and close to random on LRU traces
($13\%$)---where the optimal policy is essentially recency-driven, exactly
the pattern the model failed to acquire.

\paragraph{Why LRU is under-learnt.} Training signal is drawn from
full-cache miss events only (\emph{i.e.} the steps where Belady evicts).
Eviction density differs substantially by family:
LRU traces generate $3{,}090$ events per trace, LFU $5{,}254$, ARC $4{,}548$.
Across the 800 training traces per family this yields roughly
$2.47$M / $4.20$M / $3.64$M training examples for LRU/LFU/ARC respectively,
so LRU contributes only $\approx\!24\%$ of the gradient signal. The model
preferentially fits the heavier modes of the data distribution, which in
this dataset are LFU- and ARC-style workloads. A class-balanced sampler or
a per-family loss reweighting would likely narrow the accuracy gap on
LRU-dominant traces; this is a candidate follow-up experiment.

\paragraph{Caveat on comparing to ARC.} Unlike LRU and LFU whose eviction
rule can be computed from the cache contents alone plus access history,
ARC's decision depends on its internal T1/T2/B1/B2 partition, which is
trajectory-dependent. Our agreement metric therefore measures item-level
coincidence at shared eviction timesteps rather than slot-level agreement on
a common cache. The high coverage ($\approx\!100\%$) reflects the fact that
after warm-up, both Belady and ARC evict on every full-cache miss; the
disagreement is genuinely about which item to evict given the two
algorithms' divergent cache contents.

\subsection{Closed-loop (rollout) hit-rate evaluation}

In the teacher-forced setting the model sees Belady's cache at every event;
an eviction mistake has no downstream effect. The rollout evaluation drops
that assumption. The transformer starts each trace with an empty cache and
\emph{commits} its eviction decision on every full-cache miss; the cache
state evolves under the model's own policy, and errors cascade. Hit rate is
computed after a $10\%$ warm-up prefix (identical to the convention used at
dataset-generation time), giving one number per trace. Baseline hit rates for
LRU/LFU/ARC are taken from the recorded \texttt{*\_scores.json}, which were
computed under the same warm-up.

All 300 test traces are rolled out in parallel (batch of 300 on a single
A6000), running the model only on the subset of traces in a full-cache-miss
state at each timestep. Total rollout wall-clock: $635.6$ seconds.

\paragraph{Overall.}
\begin{center}
\begin{tabular}{lr}
\toprule
Policy & Hit rate \\
\midrule
ARC   & $62.63\%$ \\
LRU   & $57.70\%$ \\
\textbf{Model} & $\mathbf{45.61\%}$ \\
LFU   & $31.67\%$ \\
\bottomrule
\end{tabular}
\end{center}

\paragraph{Per dominating family.} Test traces are grouped by their
dominating algorithm (recorded per-trace at dataset generation time). Bold
entries mark the best policy within each row.

\begin{center}
\begin{tabular}{lrrrrr}
\toprule
Group & $n$ & Model & LRU & LFU & ARC \\
\midrule
LRU-dom & 100 & $30.55\%$ & $\mathbf{73.16\%}$ & $\phantom{0}8.81\%$ & $68.84\%$ \\
LFU-dom & 100 & $55.33\%$ & $46.54\%$ & $\mathbf{58.98\%}$ & $54.71\%$ \\
ARC-dom & 100 & $50.96\%$ & $53.40\%$ & $27.21\%$ & $\mathbf{64.33\%}$ \\
\bottomrule
\end{tabular}
\end{center}

\paragraph{Takeaways.} The closed-loop results mirror the teacher-forced
pattern almost exactly.

On \emph{LFU-dominant} traces the model is within $3.7$ pp of the LFU
baseline ($55.3\%$ vs.\ $59.0\%$), beating both LRU and ARC. This is the
clearest evidence that the learned policy is operationally useful on
workloads where frequency is the right signal.

On \emph{ARC-dominant} traces the model lags ARC by $13.4$ pp ($51.0\%$ vs.\
$64.3\%$) and is essentially tied with LRU ($53.4\%$). The model captures
some of ARC's behaviour but cannot reproduce its frequent-core-plus-drift
handling.

On \emph{LRU-dominant} traces the model collapses: $30.6\%$ hit rate against
LRU's $73.2\%$---a $42.6$ pp gap. In fact the model is closer in behaviour
to LFU ($8.8\%$) than to LRU; it has essentially no ability to track recency
in the drifting-working-set regime these traces exercise.

\paragraph{Is the model useful as a general-purpose cache policy?} No, not
as currently trained. Averaged across all three families, the model
($45.6\%$) is $12$ pp behind LRU and $17$ pp behind ARC. Even on per-family
bases it never beats the matching classical baseline. The teacher-forced
accuracy of $52\%$ might have suggested a more adaptive policy; the rollout
result clarifies that the policy is essentially LFU-shaped, and is only
competitive on LFU-shaped workloads.

\paragraph{Connection to the data-imbalance hypothesis.} The
per-dominating-family gap structure is consistent with (though not proof of)
the hypothesis in the previous subsection: LRU traces contribute only
$\approx\!24\%$ of the training events and the model correspondingly
under-learns recency tracking. Testing this directly would require at
minimum a per-family loss reweighting run; that experiment is not yet
performed.

\section{All-timesteps labelling experiment}

\subsection{Motivation and setup}

In the event-mode training above, LRU traces yielded $\approx\!24\%$ of the
training signal because they have fewer full-cache misses per trace than the
other two families. To remove this per-family imbalance we change the
labelling rule: instead of supervising the model only at actual eviction
events, we supervise at \emph{every} full-cache timestep using Belady's
\emph{hypothetical} eviction choice---the slot whose item has the furthest
next occurrence, computed regardless of whether step $t$ is a miss or a hit.
At a hit step this is guaranteed to be a slot other than the one holding
$\mathrm{trace}[t]$, since that slot's next occurrence is $t$ itself (the
minimum).

Every trace now contributes the same number of full-cache timesteps, so the
per-family gradient share becomes proportional to the number of training
traces per family ($800{:}800{:}800$, i.e.\ exactly $1{:}1{:}1$) rather than
to eviction density. Keeping the rest of the training pipeline identical
(dual-attention transformer, $d_{\text{model}}{=}128$, $d_{\text{ff}}{=}256$,
$2$ layers, $2$ heads, $w{=}1024$, $\mathrm{lr}{=}10^{-3}$), we subsample to
$4{,}000$ timesteps per trace (deterministic per-trace seed) so the training
set size matches the event-mode baseline at $9.6$M examples. Training ran
for $3$ epochs on 2${\times}$A6000 GPUs via DDP (effective batch $512$).

\subsection{Training dynamics}

Validation metrics across the three epochs (saved in
\texttt{runs/all\_timesteps/val\_loss.jsonl}):

\begin{center}
\begin{tabular}{lrrr}
\toprule
Epoch & val loss & val accuracy (all timesteps) & time \\
\midrule
0 & $1.8742$ & $49.56\%$ & $56.5$ min \\
1 & $1.8430$ & $49.82\%$ & $56.4$ min \\
2 & $1.8410$ & $49.90\%$ & $56.3$ min \\
\bottomrule
\end{tabular}
\end{center}

The smoothed training curve (\texttt{runs/all\_timesteps/training\_curves.png})
descends sharply from chance to $\approx\!2.0$ within the first $1{,}500$
steps, then plateaus near $1.85$. Train and val tracks stay on top of each
other throughout---no overfitting gap, but per-epoch improvement is tiny
($+0.3$ pp val accuracy across $3$ epochs).

Note that the validation metric here is computed over \emph{all} full-cache
timesteps (hit and miss together), which is a different distribution from
the event-mode validation above, so the values are \emph{not} directly
comparable across the two experiments.

\subsection{Closed-loop (rollout) hit-rate evaluation}

Evaluation protocol is identical to the previous subsection: start each test
trace with an empty cache, let the model commit its eviction decision on
every full-cache miss, measure hit rate after a $10\%$ warm-up.

\paragraph{Overall.}

\begin{center}
\begin{tabular}{lr}
\toprule
Policy & Hit rate \\
\midrule
ARC                      & $62.63\%$ \\
LRU                      & $57.70\%$ \\
Model (event-mode run)   & $45.61\%$ \\
\textbf{Model (all-timesteps run)} & $\mathbf{35.58\%}$ \\
LFU                      & $31.67\%$ \\
\bottomrule
\end{tabular}
\end{center}

\paragraph{Per-dominating-family comparison.}

\begin{center}
\begin{tabular}{lrrrr}
\toprule
Group & Model (event) & Model (all-ts) & $\Delta$ & Best baseline \\
\midrule
LRU-dom & $30.55\%$ & $12.93\%$ & $-17.6$ pp & LRU $73.16\%$ \\
LFU-dom & $55.33\%$ & $55.12\%$ & $-0.2$ pp  & LFU $58.98\%$ \\
ARC-dom & $50.96\%$ & $38.68\%$ & $-12.3$ pp & ARC $64.33\%$ \\
\bottomrule
\end{tabular}
\end{center}

\paragraph{Interpretation.}

Despite per-family balancing of the training distribution, the
all-timesteps model performs \emph{worse} in rollout than the event-mode
baseline across the board. The LFU-dominant regime is essentially unchanged
($-0.2$ pp); the LRU- and ARC-dominant regimes degrade substantially ($-17.6$
and $-12.3$ pp respectively).

Two effects plausibly combine to produce this outcome, neither of which we
have isolated experimentally:

\begin{itemize}
    \item \emph{Fewer actual-eviction examples.} Under the new labelling,
    each trace contributes $4{,}000$ randomly sampled full-cache timesteps,
    of which only $\approx\!27\%$ are miss-full events
    ($\approx\!1{,}080$ per trace). Across $2{,}400$ training traces this is
    $\approx\!2.6$M direct eviction examples, versus $10.3$M in the
    event-mode run---a $4\times$ reduction in the specific signal most
    relevant to the rollout metric.

    \item \emph{Hit-step supervision teaches a different skill.} At hit
    timesteps the ``which slot is worst to keep'' question is well-defined
    but operationally different from ``which slot to evict now.'' The model
    may be learning an average ranking of cached items that happens to be
    LFU-flavoured, at the expense of the specific recency heuristic that
    actual evictions demand.
\end{itemize}

Balancing the training distribution did not buy us the expected
improvement on LRU-dominant traces; rather the opposite happened. A
cleaner test of the imbalance hypothesis would be event-mode training with
explicit per-family loss reweighting, which does not dilute the eviction
signal with hit-step supervision. That experiment is not yet performed.

\subsection{Teacher-forced test evaluation}

Evaluation setup: the model sees Belady's cache state at every full-cache
timestep of the 300 held-out test traces ($4{,}783{,}043$ examples total).
Labels are Belady's hypothetical furthest-future choice. Slot-level
agreement with LRU and LFU is computed from Belady's cache at each
timestep---these rules are deterministic functions of the
(cache, access-history) pair, so coverage is trivially $100\%$. ARC is
intentionally excluded from the agreement comparison because its decision
depends on a T1/T2/B1/B2 partition that is not a function of Belady's
cache alone (see caveat in the event-mode section).

\paragraph{Full breakdown.}

\begin{center}
\small
\begin{tabular}{llrrrrr}
\toprule
Group & Subset & $n$ & Loss & Acc & agr\_LRU & agr\_LFU \\
\midrule
\multirow{3}{*}{overall} & all   & $4{,}783{,}043$ & $1.831$ & $50.03\%$ & $\phantom{0}2.23\%$ & $12.82\%$ \\
                         & event & $1{,}281{,}019$ & $1.752$ & $52.23\%$ & $\phantom{0}1.49\%$ & $12.01\%$ \\
                         & hit   & $3{,}502{,}024$ & $1.860$ & $49.23\%$ & $\phantom{0}2.50\%$ & $13.11\%$ \\
\midrule
\multirow{3}{*}{LRU-dom} & all   & $1{,}593{,}023$ & $2.774$ & $20.39\%$ & $\phantom{0}6.68\%$ & $\phantom{0}5.88\%$ \\
                         & event & $\phantom{0,}306{,}417$ & $2.897$ & $12.42\%$ & $\phantom{0}6.21\%$ & $\phantom{0}3.94\%$ \\
                         & hit   & $1{,}286{,}606$ & $2.745$ & $22.29\%$ & $\phantom{0}6.79\%$ & $\phantom{0}6.35\%$ \\
\midrule
\multirow{3}{*}{LFU-dom} & all   & $1{,}595{,}186$ & $1.447$ & $61.38\%$ & $\phantom{0}0.00\%$ & $17.60\%$ \\
                         & event & $\phantom{0,}528{,}821$ & $1.459$ & $61.20\%$ & $\phantom{0}0.00\%$ & $17.28\%$ \\
                         & hit   & $1{,}066{,}365$ & $1.441$ & $61.47\%$ & $\phantom{0}0.00\%$ & $17.76\%$ \\
\midrule
\multirow{3}{*}{ARC-dom} & all   & $1{,}594{,}834$ & $1.273$ & $68.28\%$ & $\phantom{0}0.02\%$ & $14.96\%$ \\
                         & event & $\phantom{0,}445{,}781$ & $1.313$ & $68.95\%$ & $\phantom{0}0.01\%$ & $11.32\%$ \\
                         & hit   & $1{,}149{,}053$ & $1.258$ & $68.02\%$ & $\phantom{0}0.02\%$ & $16.37\%$ \\
\bottomrule
\end{tabular}
\end{center}

\paragraph{Key observation: event-subset accuracy is unchanged.} Restricting
the all-timesteps model's accuracy to the miss-full event subset and
comparing against the event-mode model trained on those same events yields
nearly identical numbers:

\begin{center}
\begin{tabular}{lrr}
\toprule
Group (event subset) & All-timesteps model & Event-mode model \\
\midrule
Overall & $52.23\%$ & $52.46\%$ \\
LRU-dom & $12.42\%$ & $13.11\%$ \\
LFU-dom & $61.20\%$ & $61.38\%$ \\
ARC-dom & $68.95\%$ & $68.93\%$ \\
\bottomrule
\end{tabular}
\end{center}

Differences are within $0.7$ pp everywhere. \textbf{The two models produce
effectively the same eviction decisions on the same cache states.} The
$10$ pp rollout-hit-rate gap between them therefore cannot be attributed to
per-event decision quality; it must arise from cascading divergence during
closed-loop execution---one incorrect (but very-close-to-correct) eviction
choice puts the cache in a state the model is slightly worse at predicting
from, compounding over thousands of timesteps.

\paragraph{Comparing agreement metrics with the old run.} Under the
same-cache metric used here, LFU agreement peaks at $17.8\%$ on the LFU-dom
hit subset---well above the $1/32=3.1\%$ random baseline, but dramatically
lower than the $65\%$ item-level agreement we previously reported under the
(flawed) shared-timestep metric. The earlier ``the model is
LFU-flavoured'' narrative was partly an artefact of comparing against
LFU's own rollout cache rather than against what LFU would pick from
Belady's cache. With the corrected metric the model does lean towards
LFU-like decisions (never LRU-like), but the effect is much milder than
before.

\paragraph{Hit vs.\ event asymmetry.} On LRU-dominant traces, hit-step
accuracy ($22.3\%$) is \emph{nearly twice} the event-step accuracy
($12.4\%$). On LFU- and ARC-dominant traces the two are within $1$ pp.
The LRU-family hit steps are apparently easier because trace$[t]$ is in
the working set and the model only needs to avoid evicting it---picking any
non-current-item slot is usually fine, as the correct eviction is often
outside the recent-window focus. Actual miss-full events in the LRU family
require selecting the right stale phase-set item, which the model has not
learnt.

\paragraph{Summary.} The all-timesteps training matches the event-mode
model's teacher-forced accuracy exactly on real eviction events, adds some
new capability on hit steps (no direct baseline to compare against), and
loses ${\sim}10$ pp in rollout---a pure closed-loop divergence effect, not a
per-event decision issue. Fixing the LRU-dominant failure would require a
different intervention than rebalancing the labelling distribution.

\end{document}

